% =====================================================================
%  5.4 Cosinussetningen — ekstra utfordring
% =====================================================================

\begin{infoboks}{Velg metode selv}
Arket bygger på 5.4 Cosinussetningen. Tegn en figur først. Velg selv hvilken setning du bruker.
Vis regning og begrunnelse. Kontroller med vinkelsummen eller ved å sette inn.
\end{infoboks}

\nivaahode{2}

\oppg{U1}{2}{\uh}
Tre trekanter har disse sidene.
\begin{tasks}(1)
  \task Avgjør om hver trekant er spiss, rett eller stump uten å finne vinklene: $6, 7, 10$ \ og \ $8, 9, 12$ \ og \ $20, 21, 29$.
  \task Forklar hvorfor det holder å se på den lengste siden.
\end{tasks}

\oppg{U2}{2}{\mh}
En tomt $ABCD$ har sidene $AB = 35$ m, $BC = 48$ m, $CD = 42$ m og $DA = 52$ m. Diagonalen $AC$ er $60$ m.
\begin{tasks}(1)
  \task Finn vinklene $B$ og $D$.
  \task Finn arealet av tomten. Svar i hele kvadratmeter.
\end{tasks}

\oppg{U3}{2}{\prog}
Programmet skal ta tre sider og avgjøre om trekanten er spiss, rett eller stump.
\begin{kode}
from math import acos, degrees

def vinkel(a, b, c):
    """Vinkelen overfor siden a."""
    return degrees(acos((b**2 + c**2 - a**2) / (2 * b * c)))

a = 5
b = 7
c = 8

if a + b <= c or a + c <= b or b + c <= a:
    print("Ingen trekant")
else:
    A = vinkel(a, b, c)
    B = vinkel(b, a, c)
    C = vinkel(c, a, b)
    print(round(A, 1), round(B, 1), round(C, 1))
    storst = max(A, B, C)
    if ____________:
        print("rett")
    elif storst < 90:
        print("spiss")
    else:
        print("stump")
\end{kode}
\begin{tasks}(1)
  \task Fyll inn betingelsen for en rett vinkel. Hvorfor er \lstinline[style=laereverk,numbers=none,frame=none]|storst == 90| lite trygt?
  \task Test programmet på $3, 4, 5$ og $5, 7, 8$ og $1, 2, 4$. Skriv ned hva som skrives ut.
  \task Hvorfor sjekker programmet trekantulikheten først?
\end{tasks}

\nivaahode{3}

\oppg{U4}{3}{\uh}
Et parallellogram har sidene $a$ og $b$. Vinkelen mellom dem er $v$. Vis at summen av kvadratene av de to diagonalene er lik summen av kvadratene av de fire sidene.\par\smallskip
\begin{tikzpicture}[scale=0.8,font=\small]
  \coordinate (A) at (0,0); \coordinate (B) at (3.5,0);
  \coordinate (C) at (4.7,2); \coordinate (D) at (1.2,2);
  \draw[thick] (A) -- (B) node[midway, below, sloped] {$a$}
                   -- (C)
                   -- (D)
                   -- (A) node[midway, above, sloped] {$b$};
  \draw[dashed] (A) -- (C); \draw[dashed] (B) -- (D);
  \draw pic[draw, angle radius=6mm] {angle = B--A--D};
  \node at (0.95,0.3) {\scriptsize $v$};
  \node[below left] at (A) {$A$}; \node[below right] at (B) {$B$};
  \node[above right] at (C) {$C$}; \node[above left] at (D) {$D$};
\end{tikzpicture}
\hint{Hvilke trekanter har diagonalene som en av sidene? Hva er vinkelen ved $B$?}

\oppg{U5}{3}{\uh}
Lag selv en trekant med heltallige sider og en vinkel på nøyaktig $120^\circ$. Sidene skal ikke være multipler av $3, 5, 7$.
Forklar hvordan du fant den, og kontroller med cosinussetningen.

\begin{oppsummering}
Cosinussetningen $a^2 = b^2 + c^2 - 2bc\cos A$ gjelder for alle trekanter og inneholder Pytagoras som spesialtilfelle.
Sjekk først hva som er kjent, og velg setning etter det. Kontroller svaret med vinkelsum og største side mot største vinkel.
\end{oppsummering}
